\subsubsection*{NOR}

% BEGIN GENERATED ARTIFACT: def:nor-connective-propositional-logic
\begin{definitionbox}{Definition (NOR Connective)}
\begin{definition}[NOR Connective]
\label{def:nor-connective-propositional-logic}
The \emph{NOR} connective, written \(\downarrow\), is the binary connective
defined by
\[
\varphi\downarrow\psi\equiv\neg(\varphi\lor\psi).
\]
\end{definition}
\end{definitionbox}

\begin{remark*}[Standard quantified statement]
\[
\widehat v(\varphi\downarrow\psi)=\True
\Longleftrightarrow
\widehat v(\varphi)=\False\text{ and }\widehat v(\psi)=\False.
\]
\end{remark*}

\begin{remark*}[Predicate reading]
\((\varphi,\psi)\in\mathsf{NOR}\) means the negation of the disjunction of
\(\varphi\) and \(\psi\).
\end{remark*}

\begin{remark*}[Interpretation]
NOR is read as "neither." It is true exactly when both inputs are false. Like
NAND, this single connective can reproduce the standard functionally complete
basis.
\end{remark*}

\begin{dependencies}
\begin{itemize}
  \item \hyperref[def:logical-connective]{Logical Connective}
  \item \hyperref[def:truth-assignment-propositional-logic]{Truth Assignment}
\end{itemize}
\end{dependencies}
% END GENERATED ARTIFACT: def:nor-connective-propositional-logic

% BEGIN GENERATED ARTIFACT: def:nor-only-formula-propositional-logic
\begin{definitionbox}{Definition (NOR-Only Formula)}
\begin{definition}[NOR-Only Formula]
\label{def:nor-only-formula-propositional-logic}
A \emph{NOR-only formula} is a well-formed formula whose only connective is
\(\downarrow\).
\end{definition}
\end{definitionbox}

\begin{remark*}[Standard quantified statement]
\[
\varphi\text{ is NOR-only}
\Longleftrightarrow
\varphi\in\WFF_{\{\downarrow\}}.
\]
\end{remark*}

\begin{remark*}[Predicate reading]
\(\operatorname{NOROnly}(\varphi)\) means that every connective occurrence in
\(\varphi\) is an occurrence of \(\downarrow\).
\end{remark*}

\begin{remark*}[Interpretation]
A NOR-only formula uses a single primitive connective. Its expressive power
comes from the ability of NOR to define negation and disjunction.
\end{remark*}

\begin{dependencies}
\begin{itemize}
  \item \hyperref[def:nor-connective-propositional-logic]{NOR Connective}
  \item \hyperref[def:well-formed-formula]{Well-Formed Formula}
\end{itemize}
\end{dependencies}
% END GENERATED ARTIFACT: def:nor-only-formula-propositional-logic

% BEGIN GENERATED ARTIFACT: thm:nor-functionally-complete-propositional-logic
\begin{theorembox}{Theorem (NOR is Functionally Complete)}
\begin{theorem}[NOR is Functionally Complete]
\label{thm:nor-functionally-complete-propositional-logic}
The one-connective basis \(\{\downarrow\}\) is functionally complete.

\noindent\hyperref[prf:nor-functionally-complete-propositional-logic]{Go to proof.}
\end{theorem}
\end{theorembox}

\begin{remark*}[Standard quantified statement]
\[
\operatorname{FunctionallyComplete}(\{\downarrow\}).
\]
\end{remark*}

\begin{remark*}[Predicate reading]
Every Boolean function can be represented by a NOR-only formula.
\end{remark*}

\begin{remark*}[Interpretation]
NOR can define negation and disjunction:
\[
\neg\varphi\equiv\varphi\downarrow\varphi,
\qquad
\varphi\lor\psi\equiv(\varphi\downarrow\psi)\downarrow(\varphi\downarrow\psi).
\]
Once negation and disjunction are definable, functional completeness follows
from the functional completeness of \(\{\neg,\lor\}\).
\end{remark*}

\begin{dependencies}
\begin{itemize}
  \item \hyperref[def:nor-connective-propositional-logic]{NOR Connective}
  \item \hyperref[def:nor-only-formula-propositional-logic]{NOR-Only Formula}
  \item \hyperref[def:functionally-complete-connective-set-propositional-logic]{Functionally Complete Connective Set}
  \item \hyperref[thm:negation-disjunction-functionally-complete-propositional-logic]{Negation and Disjunction are Functionally Complete}
\end{itemize}
\end{dependencies}
% END GENERATED ARTIFACT: thm:nor-functionally-complete-propositional-logic
